
\documentclass[11pt]{article}
\usepackage[utf8]{inputenc}
\usepackage[T1]{fontenc}
\usepackage[italian]{babel}
\usepackage{geometry}
\geometry{a4paper, margin=2.2cm}
\usepackage{graphicx}
\usepackage{booktabs}
\usepackage{amsmath}
\graphicspath{{figures/}}

\title{Homework 2}
\author{Flavio Massaroni --- Matricola: 1990975}
\date{}

\begin{document}
\maketitle

\section{Setup and Images acquisition}

The scene is a kitchen table with an Italian Moka pot, two steel cups, one regular glass, some bananas, and a plastic water bottle. 
The images were acquired from all possible directions, at three different heights and at varying distances. 
The table is surrounded by furniture and by a plain, uniform white tiled wall next to a dark door; 
these low-textured elements were included on purpose for the geometric-artifact experiment. 
A total of 41 images were taken with fixed exposure, and copied to the computer uncompressed. Camera poses were obtained with COLMAP, 
run once on the full set of images. A fixed test set of 5 images, spread across the capture, was set aside and never used for training; 
it is identical in all experiments. The remaining 36 images form the training set of the baseline.


\section{Training Setup}
All experiments use nerfstudio with the \texttt{splatfacto} method (3D Gaussian Splatting), 30000 optimization steps, downscale factor 2, 
spherical harmonics degree 3  (default), and seed 42. No other parameter was modified, so that differences in the results can be attributed solely to the number of training views. As a metric we use PSNR:
\begin{equation*}
\mathrm{MSE} = \frac{1}{3N}\sum_{\mathbf{p}=1}^{N}\sum_{c=1}^{3}\left(\hat{I}_c(\mathbf{p}) - I_c(\mathbf{p})\right)^2,
\qquad
\mathrm{PSNR} = 10\log_{10}\frac{1}{\mathrm{MSE}}\ \mathrm{dB},
\end{equation*}
computed on the image and then mediated on the set.


\section{Baseline}
The baseline model was trained on the 36 training images. PSNR values are reported in Table~\ref{tab:psnr}.
\begin{table}[h]
\centering
\begin{tabular}{lcc}
\toprule
Run & PSNR train (dB) & PSNR test (dB) \
\midrule
Baseline (36 views) & XX.X & XX.X \
3/4 of the views (27) & XX.X & XX.X \
1/2 of the views (18) & XX.X & XX.X \
\bottomrule
\end{tabular}
\caption{Average PSNR on the training set and on the test set for each run.}
\label{tab:psnr}
\end{table}


PSNR on training views is systematically higher than on the test set: this is expected, since the radiance field is optimized to reproduce exactly those images, 
so the error on them is minimal by construction. Test views, even though close to the training ones, require a genuine interpolation of the field and 
are affected by all the flaws of the reconstruction (imprecise geometry, reflections and poorly modelled glossy surfaces). 
Figure~\ref{fig:baseline} shows two test views rendered by the baseline model, which appear photorealistic and stable.


\section{How many views are enough?}
Two additional models were trained by reducing the training set to 3/4 and 1/2 of the views, 
keeping the test set unchanged (Table~\ref{tab:psnr}). Reducing the views, the test PSNR drops progressively, 
because the field is constrained by fewer viewing rays and the Gaussians can be placed freely in the unobserved space. 
Figure~\ref{fig:views} shows the same test view rendered by the three models: with half of the views, objects lose detail and the low-textured surfaces 
(wall, door) become visibly smeared, while the central area of the table, rich in details and well covered in all runs, stays similar. 
This suggests that the number of views required is not uniform across the scene: it depends on the amount of texture and on the angular coverage of each region.

\section{A view between two cameras and outside coverage}
Two artificial views were rendered: one placed halfway between two training cameras (Figure~\ref{fig:between}) 
and one located outside the region covered by the capture, behind the scene (Figure~\ref{fig:outside}). 
The view between two cameras is almost indistinguishable from a real photo: the new position lies inside the observed region and the field is strongly 
constrained by nearby rays. The out-of-coverage view instead shows evident degradation: objects deform, surfaces "stretch", and new details are hallucinated. 
Here the radiance field is no longer a 3D model of the world but an interpolation that breaks down as soon as one leaves the hull of the training views. 
This confirms that the performance of a radiance field is reliable only inside the region covered by the photos.

\section{Geometric distribution of the radiance field}
A radiance field is optimized to reproduce the training images, not to reconstruct the scene geometry. 
Moving the camera, a clear artifact is observed above the door and on the tiled wall (Figure~\ref{fig:artifact}): in RGB, the rendering near the table 
is plausible, but the wall is fragmented into blobs, filament-like spikes radiate from the top of the door, and in the depth map (turbo colormap) 
the door does not end on a surface but the depth increases continuously upward. The reason is that these regions are low-textured 
(uniform white tiles, dark door) and appear in the training images only at the border of the frame, looking toward the table. 
The optimizer minimizes the photometric error only on the training views, and in such regions many depth configurations produce the same pixels: 
nothing penalizes Gaussians that stretch or drift along the viewing rays, until a new camera position reveals it. In other words, 
the lack of geometric constraints in these areas makes the reconstruction photometrically correct but geometrically wrong.

\section{Conclusioni}
Gli esperimenti mostrano che un radiance field va trattato come un interpolatore di
immagini vincolato dalla copertura delle viste: e affidabile tra le camere di training,
degrada allontanandosi da esse, ed e fotometricamente persuasivo anche dove la geometria e
sbagliata (superfici poco texturizzate o male coperte). Ridurre le viste peggiora le
prestazioni in modo non uniforme, colpendo prima le regioni povere di texture.

\begin{thebibliography}{9}
\bibitem{hw} Homework 2: Novel View Synthesis, Foundations of Spatial AI, Sapienza University of Rome, 2026.
\bibitem{nerfstudio} Nerfstudio: A Framework for Neural Radiance Field Development, \texttt{docs.nerf.studio}.
\end{thebibliography}

\end{document}
